\documentclass[11pt]{article}
\usepackage[margin=1in]{geometry}
\usepackage{amsmath,amssymb,amsthm}
\usepackage{hyperref}

\newtheorem{theorem}{Theorem}
\newtheorem{corollary}{Corollary}

\newcommand{\K}{\mathbb K}
\newcommand{\conv}{\operatorname{conv}}
\newcommand{\Spear}{\operatorname{Spear}}

\title{Finite-rank aDP operators attain their norm}
\author{Run fa\_banach\_001, lane 19}
\date{2026-06-16}

\begin{document}
\maketitle

\section*{Source and target}

The source paper is Kadets, Mart\'in, Mer\'i, and P\'erez,
\emph{Spear operators between Banach spaces}, arXiv:1701.02977.
All Banach spaces below are real or complex.  The source paper's
Chapter 10 asks, among other things, whether the rank-one results for
spear/aDP/lush operators extend to finite-rank operators, and also asks
whether every spear operator attains its norm.

This note proves a partial finite-rank norm-attainment result.  It does
not solve the full finite-rank lushness problem in the original codomain.

\section*{Definitions used}

A norm-one operator \(G:X\to Y\) has the alternative Daugavet property
(aDP) if
\[
  \|G+\omega T\|=1+\|T\|
\]
for every rank-one operator \(T:X\to Y\) and some unimodular scalar
\(\omega\) depending on \(T\).  A spear operator satisfies the analogous
identity for all bounded linear \(T:X\to Y\), so every spear operator
has the aDP.

We use three stated facts from the source paper.

\begin{enumerate}
\item Remark 3.3(ii): if \(G:X\to Y\) has the aDP and \(Z\subset Y\)
contains \(G(X)\), then \(G:X\to Z\) has the aDP.
\item Proposition 6.5: if the codomain is finite-dimensional, then for
a norm-one operator \(H:X\to Z\), lushness, spearness, and the aDP are
equivalent.
\item Proposition 7.9: if \(H:X\to Z\) is lush, then
\[
  B_X=\overline{\conv}\{x\in S_X:\|Hx\|=1\}.
\]
\end{enumerate}

\section*{Proof Intuition}

The finite-rank assumption lets us replace the original codomain by the
actual range \(Z=G(X)\).  That move is one-way for the aDP: the source
paper explicitly allows codomain restriction, while also warning that
codomain extension can fail.  Once the codomain is the finite-dimensional
space \(Z\), the source paper's finite-dimensional codomain theorem
upgrades the restricted operator from aDP to lush.  Lushness then gives
norm-attainment, and because \(Z\) carries the norm inherited from \(Y\),
the norm-attaining vectors for the restricted operator are exactly the
norm-attaining vectors for the original finite-rank operator.

\section*{Main result}

\begin{theorem}
Let \(X,Y\) be Banach spaces and let \(G:X\to Y\) be a norm-one
finite-rank operator.  If \(G\) has the alternative Daugavet property,
then
\[
  B_X=\overline{\conv}\{x\in S_X:\|Gx\|=1\}.
\]
In particular, \(G\) attains its norm.
\end{theorem}

\begin{proof}
Let \(Z=G(X)\), endowed with the norm inherited from \(Y\), and let
\(H:X\to Z\) be the same algebraic operator, \(Hx=Gx\).  Since \(G\)
has finite rank, \(Z\) is finite-dimensional.  Also \(\|H\|=\|G\|=1\).

By Remark 3.3(ii) of the source paper, the aDP of \(G:X\to Y\) passes
to the restricted-codomain operator \(H:X\to Z\).  Since \(Z\) is
finite-dimensional, Proposition 6.5 of the source paper says that
\(H\) is lush.  Proposition 7.9 of the same paper therefore yields
\[
  B_X=\overline{\conv}\{x\in S_X:\|Hx\|=1\}.
\]
Finally, \(Z\) has the norm inherited from \(Y\), so \(\|Hx\|=\|Gx\|\)
for every \(x\in X\).  This is exactly the displayed conclusion.
\end{proof}

\begin{corollary}
Every finite-rank spear operator \(G:X\to Y\) attains its norm.  More
precisely,
\[
  B_X=\overline{\conv}\{x\in S_X:\|Gx\|=\|G\|\}
\]
for every nonzero finite-rank spear operator \(G\).
\end{corollary}

\begin{proof}
After normalizing \(G\), spearness implies the aDP, so the theorem
applies.  Scaling back gives the displayed equality.
\end{proof}

\section*{What remains open}

The proof deliberately avoids claiming that \(G:X\to Y\) is lush in the
original codomain.  The source paper states that aDP is preserved when
one restricts the codomain to a subspace containing the range, but not
under arbitrary codomain extension.  Thus the range-restricted operator
\(H:X\to G(X)\) is lush, while the original finite-rank operator
\(G:X\to Y\) may still resist the rank-one-to-finite-rank lushness
problem.  The same obstruction prevents the argument from proving the
full finite-rank analogue of the source paper's adjoint-lushness
rank-one result.

\section*{Verification and novelty notes}

The proof is a formal combination of three labeled results in
arXiv:1701.02977: Remark 3.3(ii), Proposition 6.5, and Proposition 7.9.
The set of norm-attaining vectors is unchanged under the range-codomain
restriction because the range carries the inherited norm.

Cheap run-index searches found no prior packet for arXiv:1701.02977 or
for this finite-rank aDP norm-attainment consequence.  Local checks
against the later papers arXiv:2306.02645 and arXiv:2306.02638 did not
find this statement.  Web searches on 2026-06-16 for exact combinations
of ``finite-rank'', ``alternative Daugavet property'', ``spear
operators'', and ``norm-attaining'' found no exact duplicate.  Because
the proof is a short derived corollary from the source paper, this
packet should be reviewed as a modest but useful partial result, not as
a full solution of the finite-rank lushness problem.

\begin{thebibliography}{9}

\bibitem{KMMP}
V. Kadets, M. Mart\'in, J. Mer\'i, and D. P\'erez,
\emph{Spear operators between Banach spaces},
arXiv:1701.02977.

\end{thebibliography}

\end{document}

